\section{Related Work}
\label{sec:related}
\paragraph{Feedforward 3D reconstruction}
methods address the problem of 3D object reconstruction by learning a feedforward model in a regression-based manner.
Earlier works~\cite{choy20163d,wang2018pixel2mesh,mescheder2019occupancy,xu2019disn,groueix2018,wu2017marrnet,huang2023shapeclipper} in this field typically predict only the geometry and train on small datasets~\cite{chang2015shapenet,sun2018pix3d}, which limits their generalization ability. Recently, larger 3D datasets~\cite{reizenstein2021co3d,deitke2022objaverse} have been collected, unlocking the potential to train feedforward 3D models at scale~\cite{wu2023multiview,huang2024zeroshape,hong2023lrm}. These models exhibit great generalization ability to unseen images, and excel at producing reconstructions that tightly align with the observed cues in the input image. In particular, LRM~\cite{hong2023lrm} and follow-up works~\cite{TripoSR2024,tang2024lgm,zhang2024gs,xu2024instantmesh,wang2024crm,sf3d2024} show that properly designed large transformer models can be trained using only rendering losses to capture object geometry and texture in great detail. Despite the high fidelity and computational efficiency of these models, the oversimplified bijective assumption in these regressive approaches results in oversmoothed unseen surfaces. Multi-view diffusion models~\cite{liu2023zero,shi2023mvdream,liu2023syncdreamer,shi2023zero123plus} have been considered as a remedy for this, where additional viewpoints are synthesized as input to the feedforward model~\cite{tang2024lgm,wang2024crm,xu2024instantmesh}. However, the inconsistency across viewpoints often leads to significant artifacts on the reconstructed surfaces, and the computational efficiency of these approaches is severely affected by the slow multi-view generation process. Our model also aims to overcome the learning ambiguity in regressive approaches, but our point sampling approach is inherently 3D-consistent and computationally efficient, and further allows easy user edits.

\begin{figure*}[t]
\centering
\includegraphics[width=1.0\linewidth]{figures/overview.pdf}
\caption{\textbf{\method Overview.} Conditioned on the input image, \method first leverages a point diffusion model to generate a sparse point cloud. The triplane transformer then uses the sampled point cloud and image features to produce high-resolution triplane features. The triplane features are then queried to reconstruct the geometry, texture, and illumination of the object in the image.}
\vspace{-5mm}
\label{fig:overview}
\end{figure*}

\paragraph{Generative 3D modeling}
learns the image-conditioned distribution of 3D assets instead of a deterministic mapping. Early 3D generative works use GAN~\cite{valsesia2018learning,hui2020progressive,cai2020learning,gao2022get3d}, normalizing flow~\cite{yang2019pointflow,klokov2020discrete} or VAE~\cite{wu2019sagnet,mittal2022autosdf,gao2021tm} as the generative framework. Inspired by the success of 2D diffusion models~\cite{dhariwal2021diffusion,rombach2022high}, 3D diffusion models~\cite{nichol2022point,jun2023shap,liu2023one2345++,shue20233d,chou2023diffusion,shim2023diffusion,cheng2023sdfusion,li2023diffusion,liu2023zero,shi2023mvdream,liu2023syncdreamer,huang2024pointinfinity,yariv2024mosaic,chen2023single,zhao2023michelangelo,lan2024ln3diff} have also been extensively explored in recent works. Despite the advantage of probabilistic modeling that avoids over-smoothed results, diffusion-based 3D generation has two drawbacks: 1) not aligning well with input observations, and 2) having low inference speed at high resolution. Our work inherits the advantage of probabilistic modeling, while avoiding the drawbacks by using diffusion to generate only sparse point clouds.

\paragraph{Optimization-based single-view 3D} 
leverages 2D generative priors to recover 3D from single-view images. These works~\cite{melas2023realfusion,deng2023nerdi,tang2023make,gu2023nerfdiff} rely on SDS-type loss~\cite{pooledreamfusion, wang2023score} and generate 3D assets by optimizing for each object image separately. These methods achieve promising results without large-scale annotation. However, the lack of a strong explicit 3D prior makes the optimization process inefficient and prone to local minima.